\documentclass[12pt]{article}
\usepackage[utf8]{inputenc}
\usepackage[T1]{fontenc}
\usepackage{hyperref}
\title{Wutunhua Language Resources for K-12 Curriculum Development}
\author{Compiled by Claude Sonnet 5.5 (claude-sonnet-5-5) for Andr\'as Kornai}
\date{}
\begin{document}
\maketitle

\section{Introduction}
Wutunhua, or Wutun (ISO 639-3 wuh, Glottocode wutu1241), is a mixed language spoken in Tongren County, Qinghai, China. Its vocabulary and much of its grammar come from Northwest Mandarin, but its morphology and syntax are strongly influenced by Amdo Tibetan and Mongolic (Bonan). Search results give about 4,000 speakers, and the reference description gives about 2,000. Speakers identify as Tibetan. Some sources call the language Ngandehua. Some call it a creole, while others call it a mixed language.

\section{Primary Educational Resources}
No primer, textbook or school programme was found, and the language is not known to be taught in schools.

\section{Linguistic and Reference Materials}
Ke Sandman's work is the main grammatical source.
\begin{itemize}
\item Wutun language (Wikipedia): \url{https://en.wikipedia.org/wiki/Wutun_language}
\item Sandman, ``Wutun as a mixed language'' (De Gruyter, 2021), abstract and details: \url{https://www.jbe-platform.com/content/journals/10.1075/jpcl.00103.san}
\item Sandman, paper in \emph{Sustainable Language} 264, on Bonan and Mongolic influence: \url{https://www.sgr.fi/sust/sust264/sust264_sandman.pdf}
\item A University of Helsinki repository document found in the search: \url{https://helda.helsinki.fi/server/api/core/bitstreams/5e2bfb82-eda7-4d95-9bf3-8def8ef92b4d/content}
\item Language Log, ``Mixed'' languages: \url{https://languagelog.ldc.upenn.edu/nll/?p=29508}
\end{itemize}
Janhunen (2007) is cited in the results for the Qinghai Linguistic Complex, but no URL was seen for it.

\section{Religious and Community Resources}
No scripture was found. A missions-oriented people-group profile of the Wutun was found. It is not a language resource.
\begin{itemize}
\item Asia Harvest, Wutun people-group profile: \url{https://people-groups.asiaharvest.org/people-group-profiles/qinghai/wutun.pdf}
\end{itemize}

\section{Digital and Audio Resources}
The Endangered Languages Documentation Programme archive holds primary data of Wutunhua: \url{https://eldp.access.preservica.com/dk0020/}. The search result did not show the contents, so the types of recording it holds are not confirmed.

\section{Curriculum Development Context}
Wutun has a strong scholarly description but no known teaching materials or orthography. A curriculum would require community-led orthography development, building on Sandman's grammar and the archived data. Mandarin and Tibetan are the schooling languages in the region.
\end{document}